% §12.4 Bridge
%
% Cross-reference resolution status:
%
% DEFINED here:
%   \label{sec:ch12-bridge}                       -- this section
%
% RESOLVED (within chapter 12, drafted in Tasks 2-4 of Plan 12):
%   \Cref{sec:sets-to-relations}                  -- §12.1
%   \Cref{sec:relation-composition}               -- §12.2
%   \Cref{thm:join-rates}                         -- §12.2
%   \Cref{sec:approximate-joins}                  -- §12.3
%
% RESOLVED (earlier chapters):
%   \Cref{ch:approximate-sets}                    -- chapter 5
%   \Cref{ch:set-composition}                     -- chapter 6
%   \Cref{ch:approximate-functions}               -- chapter 9
%   \Cref{sec:minhash-lsh}                        -- chapter 10 §10.4
%
% FORWARD REFS to later chapters (Plan 0 stubs):
%   \Cref{sec:bib-notes-ch12}                     -- §12.5 (Task 6)
%   \Cref{ch:algebra-types}                       -- chapter 13

\section{Bridge}
\label{sec:ch12-bridge}

The three prior sections of this chapter trace one object from three
sides. \Cref{sec:sets-to-relations} reads an approximate relation as
the chapter~5 framework (\Cref{ch:approximate-sets}) applied to the
pair universe $T \times U$: the construction is forced by the
classical correspondence between a relation and its membership
predicate, the two rates $(\fprate, \fnrate)$ retain their chapter~5
meaning with pairs in place of elements, and the chapter~9 binary-
function reading (\Cref{ch:approximate-functions}) is an equivalent
view of the same object. \Cref{sec:relation-composition} adds the
relational operators (natural join, projection, selection) and works
out their rate-propagation rules. The headline is
\Cref{thm:join-rates}: under cross-relation independence, two
FNR-zero approximate relations with FPRs $\fprate_1$ and $\fprate_2$
join into a relation with $\fnrate_{\bowtie} = 0$ and
$\fprate_{\bowtie} = \fprate_1 \fprate_2$, the same product structure
chapter~6 (\Cref{ch:set-composition}) derived for intersection of
approximate sets. \Cref{sec:approximate-joins} instantiates the
machinery on the MinHash + LSH similarity-join pipeline of chapter~10
(\Cref{sec:minhash-lsh}) and reads the published FPR/FNR analyses of
that pipeline as the chapter~12 join rates specialized to a
particular sketch-based predicate. Three sections, one approximate
relation: a definition, a composition algebra, an algorithmic
instance.

The bibliographic notes in \Cref{sec:bib-notes-ch12} name the
antecedents that the prose has been compressing: the
\textcite{towell-bernoulli-relations} manuscript is the primary
source for the relation-level framework drafted here, Codd's
relational algebra is the classical model that the approximate side
specializes, and the database literature on similarity joins
(MinHash-LSH and equijoin variants) is the algorithmic body of work
that the chapter reads through the framework. The notes are short
because the technical content is largely a lift: most of the work was
done in chapters~5, 6, and 10, and chapter~12 is the relation-level
specialization.

\Cref{ch:algebra-types} closes Part~III by giving the type-driven
algebra that unifies the chapter on approximate sets, the chapter on
approximate functions, the present chapter on approximate relations,
and the chapter on sketches into a single algebra of $\Bernoulli{T}$
constructors. An approximate set is $\Bernoulli{\BoolType}$ on a
universe; an approximate relation is $\Bernoulli{\BoolType}$ on a
product universe $T \times U$ (equivalently $\Bernoulli{T \times U}$
in the type-constructor view); a sketch over a list type is
$\Bernoulli{L\,T}$; a Monte Carlo decision procedure is
$\Bernoulli{\BoolType}$. Composition rules that each of the prior
chapters derived by direct calculation become functorial properties
of the type algebra. The next chapter does that unification
explicitly, and chapter~12 is one of its first instances: an
approximate relation reread as $\Bernoulli{T \times U}$ inherits its
composition rules from the algebraic structure of the product type
rather than from a separate derivation.
